\chap{DC Circuit Analysis -- First Project on LTspice}

\section{Introduction}
This first laboratory introduces LTspice, the circuit simulator used throughout this manual. You will build resistive circuits, run a DC operating-point analysis and compare the simulated values with your own hand calculations.

\textbf{After this laboratory you will be able to:}
\begin{itemize}
  \item create, wire and save a schematic in LTspice;
  \item set component values with SPICE unit prefixes and place ideal sources and dependent sources;
  \item run a DC operating-point analysis (\cmd{.op}) and read node voltages and branch currents;
  \item verify Kirchhoff's laws, equivalent resistance, circuit transformations and power balance by simulation;
  \item compare theory with simulation and explain the differences;
  \item use parameters (\cmd{.param}), parametric steps (\cmd{.step}) and measurements (\cmd{.meas}) to automate the comparison.
\end{itemize}

\textbf{How the exercises are organised.} Every exercise has the same structure: the statement and the circuit figure; a \emph{Theory calculation} part that you complete by hand; an \emph{LTspice simulation} part that tells you what to build and which analysis to run; and a \emph{Comparison} table where you put your results side by side. The calculation space is intentionally left blank. Save every schematic as \file{Lab1/ExNN.asc} (\file{Ex01.asc} to \file{Ex10.asc} for main exercises; \file{Ex11.asc} and \file{Ex12.asc} for advanced exercises).

\section{LTspice guide}
\subsection{Create a schematic}
Start LTspice and choose \textit{File $\rightarrow$ New Schematic} (Figure~\ref{fig:lab1_guide01}). Save it immediately (\textit{File $\rightarrow$ Save As}) as \file{Lab1/Ex01.asc}.
\labfig{lab1_guide01.png}{10cm}{Create a new schematic in LTspice}{fig:lab1_guide01}

\subsection{Place and wire components}
\begin{itemize}
  \item Press \texttt{R} to place a resistor, \texttt{G} to place the ground (node \texttt{0}). Press \texttt{F2} to open the component list: search \cmd{voltage} for a DC voltage source, \cmd{current} for a current source, and \cmd{e}, \cmd{f}, \cmd{g}, \cmd{h} for the dependent sources.
  \item Rotate a component with \texttt{Ctrl+R} and mirror it with \texttt{Ctrl+E} before clicking to place it.
  \item Press \texttt{F3} (or the pencil icon) to draw wires between component pins. Press \texttt{F5} to delete and \texttt{F7} to move an object.
  \item Press \texttt{F4} to give a net a name (\textit{Label Net}), for example \texttt{A}, \texttt{B}.
  \item \textbf{Every circuit needs a ground} (node \texttt{0}); otherwise the simulation cannot run.
\end{itemize}
\labfig{lab1_guide02.png}{12cm}{Placing components and wires}{fig:lab1_guide02}

\subsection{Set the component values}
Right-click a component to open its dialog and type the value; for a voltage source enter the DC value. Use SPICE prefixes: \texttt{f}, \texttt{p}, \texttt{n}, \texttt{u} (micro), \texttt{m} (milli), \texttt{k}, \texttt{Meg}, \texttt{G}. \textbf{Warning:} in SPICE \texttt{m} and \texttt{M} both mean milli; mega must be written \texttt{Meg}.
\labfig{lab1_guide03.png}{9cm}{Editing the value of a component}{fig:lab1_guide03}

\subsection{Add the simulation command and run}
Choose \textit{Simulate $\rightarrow$ Edit Simulation Cmd}, select the \textit{DC op pnt} tab and click OK; the directive \cmd{.op} is placed on the schematic (you can also place it as text with the \texttt{S} key). Then choose \textit{Simulate $\rightarrow$ Run}.
\labfig{lab1_guide04.png}{10cm}{Adding the \texttt{.op} directive}{fig:lab1_guide04}

\subsection{Read the results}
After a \cmd{.op} run, a window lists the node voltages and the branch currents (Figure~\ref{fig:lab1_guide05}). If you use \cmd{.meas}, its results are written in the SPICE Error Log (\textit{View $\rightarrow$ SPICE Error Log}). Note the sign conventions: the current of a voltage source is positive when it flows into its positive terminal, so a source that delivers power shows a negative current.
\labfig{lab1_guide05.png}{10cm}{Operating-point results}{fig:lab1_guide05}

\section{Exercises}

\subsection{Exercise 1: Kirchhoff's voltage law with a dependent source}
\status{Unchanged (direct transfer from Ch.~1 on PSpice).}
Consider the circuit in Figure~\ref{fig:lab1_ex01}. It contains a 12\,V source, a 4\,V source, a 4\,$\Omega$ resistor, a 6\,$\Omega$ resistor and a voltage-controlled voltage source of value $2v_0$ (positive terminal on the left), where $v_0$ is the voltage across the 6\,$\Omega$ resistor (positive terminal on the left). Calculate the voltage $v_0$ and the current $i$ using Kirchhoff's voltage law (KVL). Then simulate the circuit to check your results.

\labfig{lab1_ex01.png}{7cm}{Exercise 1: circuit with a dependent source}{fig:lab1_ex01}

\theory{8}

\simulation{
\item Draw the circuit of Figure~\ref{fig:lab1_ex01} in LTspice (components \cmd{res}/\cmd{voltage}/\cmd{current}, ground \cmd{0}) and save it as \file{Lab1/Ex01.asc}.
\item Use the voltage-controlled voltage source \cmd{e} (press \texttt{F2} and search for \cmd{e}). Set its gain to 2 and connect its two control pins across the 6\,$\Omega$ resistor, with the polarity of $v_0$.
\item Add the \cmd{.op} analysis (\textit{Simulate $\rightarrow$ Edit Simulation Cmd $\rightarrow$ DC op pnt}), run the simulation and read the node voltages and branch currents in the operating-point window.
\item Be careful with signs: LTspice reports the current of a voltage source as positive when it flows \emph{into} its positive terminal, and the current of a resistor as positive when it flows from its first pin to its second pin.
}

\comparetable{
\cmprow{$v_0$ (V)}
\cmprow{$i$ (A)}
}

\commentlines{2}

\subsection{Exercise 2: Equivalent resistance by series and parallel reductions}
\status{Unchanged (direct transfer from Ch.~1 on PSpice).}
Consider the circuit in Figure~\ref{fig:lab1_ex02}. Rearrange the circuit to clarify its series and/or parallel topology, and draw the rearranged circuit. Then use series/parallel reductions to find the equivalent resistance seen by the source between terminals A and F. Finally, simulate the circuit to check that the current through the whole circuit is correctly calculated.

\labfig{lab1_ex02.png}{10cm}{Exercise 2: find the equivalent resistance between terminals A and F}{fig:lab1_ex02}

\theory{8}

\simulation{
\item Draw the circuit of Figure~\ref{fig:lab1_ex02} in LTspice (components \cmd{res}/\cmd{voltage}/\cmd{current}, ground \cmd{0}) and save it as \file{Lab1/Ex02.asc}.
\item Add the \cmd{.op} analysis (\textit{Simulate $\rightarrow$ Edit Simulation Cmd $\rightarrow$ DC op pnt}), run the simulation and read the node voltages and branch currents in the operating-point window.
\item Label the nodes A to F with \textit{Label Net} (\texttt{F4}), so that the results are listed with these names.
\item The equivalent resistance is $R_{AF}=V_1/I$, where $I$ is the magnitude of the current delivered by the source.
\item Be careful with signs: LTspice reports the current of a voltage source as positive when it flows \emph{into} its positive terminal, and the current of a resistor as positive when it flows from its first pin to its second pin.
}

\comparetable{
\cmprow{$R_{AF}$ ($\Omega$)}
\cmprow{$I$ (A)}
}

\commentlines{2}

\subsection{Exercise 3: Equivalent resistance and node voltages}
\status{Unchanged (direct transfer from Ch.~1 on PSpice).}
Consider the circuit in Figure~\ref{fig:lab1_ex03}. Rearrange the circuit to clarify its series and/or parallel topology. Next, find the equivalent resistance between terminals A and F and the voltages at A, B, C, D and E. Finally, simulate the circuit to check your calculation.

\labfig{lab1_ex03.png}{11cm}{Exercise 3: equivalent resistance and node voltages}{fig:lab1_ex03}

\theory{10}

\simulation{
\item Draw the circuit of Figure~\ref{fig:lab1_ex03} in LTspice (components \cmd{res}/\cmd{voltage}/\cmd{current}, ground \cmd{0}) and save it as \file{Lab1/Ex03.asc}.
\item Add the \cmd{.op} analysis (\textit{Simulate $\rightarrow$ Edit Simulation Cmd $\rightarrow$ DC op pnt}), run the simulation and read the node voltages and branch currents in the operating-point window.
\item Label the nodes A to F (\texttt{F4}); node F is the ground. The source is a current source: place it with \cmd{current} and check the arrow direction.
\item Be careful with signs: LTspice reports the current of a voltage source as positive when it flows \emph{into} its positive terminal, and the current of a resistor as positive when it flows from its first pin to its second pin.
}

\comparetable{
\cmprow{$R_{AF}$ ($\Omega$)}
\cmprow{$V_A$ (V)}
\cmprow{$V_B$ (V)}
\cmprow{$V_C$ (V)}
\cmprow{$V_D$ (V)}
\cmprow{$V_E$ (V)}
}

\commentlines{2}

\subsection{Exercise 4: Currents and voltages in a resistive network}
\status{Unchanged (direct transfer from Ch.~1 on PSpice).}
Given the circuit in Figure~\ref{fig:lab1_ex04}, find $I_1$, $I_2$, $I_3$, $V_a$ and $V_b$. Present your calculation steps and check them by simulation.

\labfig{lab1_ex04.png}{8cm}{Exercise 4: find $I_1$, $I_2$, $I_3$, $V_a$ and $V_b$}{fig:lab1_ex04}

\theory{8}

\simulation{
\item Draw the circuit of Figure~\ref{fig:lab1_ex04} in LTspice (components \cmd{res}/\cmd{voltage}/\cmd{current}, ground \cmd{0}) and save it as \file{Lab1/Ex04.asc}.
\item Add the \cmd{.op} analysis (\textit{Simulate $\rightarrow$ Edit Simulation Cmd $\rightarrow$ DC op pnt}), run the simulation and read the node voltages and branch currents in the operating-point window.
\item Compare the simulated branch currents with $I_1$, $I_2$, $I_3$ (check their reference directions).
\item Be careful with signs: LTspice reports the current of a voltage source as positive when it flows \emph{into} its positive terminal, and the current of a resistor as positive when it flows from its first pin to its second pin.
}

\comparetable{
\cmprow{$I_1$ (mA)}
\cmprow{$I_2$ (mA)}
\cmprow{$I_3$ (mA)}
\cmprow{$V_a$ (V)}
\cmprow{$V_b$ (V)}
}

\commentlines{2}

\subsection{Exercise 5: Selecting a standard resistor}
\status{Modified (added \cmd{.param} and \cmd{.step} parametric sweep to investigate standard resistor selection).}
Given the network in Figure~\ref{fig:lab1_ex05}: a 24\,V source in series with an unknown resistor $R$ and a 3.9\,k$\Omega$ resistor, across which the voltage must be $V_1=5$\,V. Select resistor $R$ from the standard resistor list (table at the end of this lab) and complete the following requirements.

\labfig{lab1_ex05.png}{6cm}{Exercise 5: choose $R$ from the standard resistor list}{fig:lab1_ex05}

\textit{\textbf{Notes:}}\\
\textit{Explanations, formulas, and equations are expected rather than only results.}
\begin{enumerate}[label=\alph*.]
  \item Find the required value of the resistor $R$.\\ $R=$\dotfill
  \item Use the standard values of the table at the end of this lab to select a resistor with 10\,\% tolerance for $R$. $R$ may be a single resistor or a combination of resistors, as long as every resistor has a standard value and is available on the market.\\ The selected resistor:\dotfill
  \item Using the resistor selected in (b), determine the voltage across the 3.9\,k$\Omega$ resistor.\\ $V_1=$\dotfill
  \item Calculate the percent error in $V_1$ when the standard resistor of (b) is used.\\ \dotfill
  \item Determine the power rating for this standard component (choose among standard ratings such as 1/8, 1/4, 1/2 and 1\,W).\\ $P_R=$\dotfill
\end{enumerate}

\theory{0}

\simulation{
\item Draw the circuit of Figure~\ref{fig:lab1_ex05} in LTspice (components \cmd{res}/\cmd{voltage}/\cmd{current}, ground \cmd{0}) and save it as \file{Lab1/Ex05.asc}.
\item Define the unknown resistor as a parameter: place a \cmd{.param} directive (\texttt{S}) such as \cmd{.param R=<value>} and type \cmd{\{R\}} in the value field of resistor R.
\item Add the \cmd{.op} analysis (\textit{Simulate $\rightarrow$ Edit Simulation Cmd $\rightarrow$ DC op pnt}), run the simulation and read the node voltages and branch currents in the operating-point window.
\item Use \cmd{.step param R list ...} with the selected value and its two neighbours in the standard list to see how $V_1$ changes (each step is one simulation run).
\item Be careful with signs: LTspice reports the current of a voltage source as positive when it flows \emph{into} its positive terminal, and the current of a resistor as positive when it flows from its first pin to its second pin.
}

\comparetable{
\cmprow{$R$ required ($\Omega$)}
\cmprow{$R$ selected ($\Omega$)}
\cmprow{$V_1$ (V)}
\cmprow{Error in $V_1$ (\%)}
\cmprow{$P_R$ (mW)}
}

\commentlines{2}

\subsection{Exercise 6: Circuit transformation (Y--$\Delta$)}
\status{Unchanged (direct transfer from Ch.~1 on PSpice).}
Consider the circuit in Figure~\ref{fig:lab1_ex06}. Apply what you have learned to transform it into another form in which the total equivalent resistance $R_{ab}$ is found more easily, and draw the transformed circuit. Next, find the current $i$ through the circuit and simulate it to check your results.

\labfig{lab1_ex06.png}{9cm}{Exercise 6: find $R_{ab}$ and the current $i$}{fig:lab1_ex06}

\theory{8}

\simulation{
\item Draw the circuit of Figure~\ref{fig:lab1_ex06} in LTspice (components \cmd{res}/\cmd{voltage}/\cmd{current}, ground \cmd{0}) and save it as \file{Lab1/Ex06.asc}.
\item Add the \cmd{.op} analysis (\textit{Simulate $\rightarrow$ Edit Simulation Cmd $\rightarrow$ DC op pnt}), run the simulation and read the node voltages and branch currents in the operating-point window.
\item $R_{ab}=V_1/i$; read $i$ from the current of the source (mind the sign).
\item Be careful with signs: LTspice reports the current of a voltage source as positive when it flows \emph{into} its positive terminal, and the current of a resistor as positive when it flows from its first pin to its second pin.
}

\comparetable{
\cmprow{$R_{ab}$ ($\Omega$)}
\cmprow{$i$ (A)}
}

\commentlines{2}

\subsection{Exercise 7: Circuit transformation and total current}
\status{Unchanged (direct transfer from Ch.~1 on PSpice).}
Consider the circuit in Figure~\ref{fig:lab1_ex07}. Transform it into another form in which the total equivalent resistance is found more easily, and draw the transformed circuit. Next, find the current $I_S$ through the circuit and simulate it to check your results.

\labfig{lab1_ex07.png}{6cm}{Exercise 7: find the equivalent resistance and $I_S$}{fig:lab1_ex07}

\theory{8}

\simulation{
\item Draw the circuit of Figure~\ref{fig:lab1_ex07} in LTspice (components \cmd{res}/\cmd{voltage}/\cmd{current}, ground \cmd{0}) and save it as \file{Lab1/Ex07.asc}.
\item Add the \cmd{.op} analysis (\textit{Simulate $\rightarrow$ Edit Simulation Cmd $\rightarrow$ DC op pnt}), run the simulation and read the node voltages and branch currents in the operating-point window.
\item Enter resistor values with the SPICE prefix \texttt{k} (for example \texttt{12k}); never use \texttt{M} for mega, because in SPICE \texttt{m} and \texttt{M} both mean milli (mega is \texttt{Meg}).
\item Be careful with signs: LTspice reports the current of a voltage source as positive when it flows \emph{into} its positive terminal, and the current of a resistor as positive when it flows from its first pin to its second pin.
}

\comparetable{
\cmprow{$R_{eq}$ (k$\Omega$)}
\cmprow{$I_S$ (mA)}
}

\commentlines{2}

\subsection{Exercise 8: Active and passive elements (1)}
\status{Modified (represented unknown elements with labelled boxes and equivalent sources/resistors).}
In the circuit of Figure~\ref{fig:lab1_ex08}, the elements $p_2$, $p_3$ and $p_4$ are unknown elements; the voltage across each element and the current through it are given in the figure, together with a 30\,V source $p_1$ and a current-controlled current source $p_5$ of value $0.4I$. First, use the passive sign convention to identify whether each of $p_2$, $p_3$, $p_4$ \emph{absorbs} or \emph{supplies} power. Represent an absorbing element by a pure resistor with a proper value, and a supplying element by an ideal DC voltage source with the corresponding value. Next, redraw the circuit and calculate the power absorbed by each element (verify that the total power is zero). Then perform a simulation with the elements determined in the previous step.

\labfig{lab1_ex08.png}{10cm}{Exercise 8: determine the unknown elements and the absorbed power of each element}{fig:lab1_ex08}

\theory{8}

\simulation{
\item Draw the original circuit with $p_2$, $p_3$, $p_4$ as labelled rectangles (\textit{Draw $\rightarrow$ Rectangle}) and add the given voltages and currents as text. This drawing is only used for the statement; it is not simulated.
\item Redraw the circuit with the resistors/voltage sources you determined and save it as \file{Lab1/Ex08.asc}.
\item The current-controlled current source is \cmd{f}: its controlling current is the current of a voltage source. To sense the controlling current $I$, insert an independent 0\,V DC voltage source in series with the branch of $p_2$, and reference its name along with the gain 0.4 in the value field of \cmd{f} (e.g., \cmd{Vsense 0.4}).
\item Add the \cmd{.op} analysis (\textit{Simulate $\rightarrow$ Edit Simulation Cmd $\rightarrow$ DC op pnt}), run the simulation and read the node voltages and branch currents in the operating-point window.
\item Compute the power of each element from the simulated voltages and currents.
\item Be careful with signs: LTspice reports the current of a voltage source as positive when it flows \emph{into} its positive terminal, and the current of a resistor as positive when it flows from its first pin to its second pin.
}

\comparetable{
\cmprow{$p_2$: type and power (W)}
\cmprow{$p_3$: type and power (W)}
\cmprow{$p_4$: type and power (W)}
\cmprow{$p_1$ (W)}
\cmprow{$p_5$ (W)}
}

\commentlines{2}

\subsection{Exercise 9: Active and passive elements (2)}
\status{Modified (represented unknown elements with labelled boxes and equivalent sources/resistors).}
Consider the circuit in Figure~\ref{fig:lab1_ex09}, where a 10\,V source supplies a 12\,$\Omega$ resistor, the elements $p_1$ and $p_2$ are unknown, and a current-controlled voltage source of value $3i_x$ closes the loop; node D is the ground. Find the voltage $v$ and the current $i_x$. According to the result, determine whether the elements whose absorbed powers are $p_1$ and $p_2$ are active or passive (calculations are required); use the passive sign convention. Represent an element that consumes power by a pure resistor with an appropriate value, and a power-supplying element by a corresponding ideal DC voltage source. Perform a simulation to check how the circuit works.

\labfig{lab1_ex09.png}{8cm}{Exercise 9: find the unknown elements and variables}{fig:lab1_ex09}

\theory{10}

\simulation{
\item Draw the original circuit with $p_1$, $p_2$ as labelled rectangles (statement figure only), then redraw the circuit with your resistor / voltage source and save it as \file{Lab1/Ex09.asc}.
\item Use the current-controlled voltage source \cmd{h}: its controlling current is the current of a voltage source. To sense $i_x$, insert an independent 0\,V DC voltage source in series with the branch carrying $i_x$, and reference its name and transresistance 3 in the value field of \cmd{h} (e.g., \cmd{Vsense 3}).
\item Add the \cmd{.op} analysis (\textit{Simulate $\rightarrow$ Edit Simulation Cmd $\rightarrow$ DC op pnt}), run the simulation and read the node voltages and branch currents in the operating-point window.
\item Be careful with signs: LTspice reports the current of a voltage source as positive when it flows \emph{into} its positive terminal, and the current of a resistor as positive when it flows from its first pin to its second pin.
}

\comparetable{
\cmprow{$v$ (V)}
\cmprow{$i_x$ (A)}
\cmprow{$p_1$: type and power (W)}
\cmprow{$p_2$: type and power (W)}
}

\commentlines{2}

\subsection{Exercise 10: Resistive network: find a voltage}
\status{Unchanged (direct transfer from Ch.~1 on PSpice).}
Given the circuit in Figure~\ref{fig:lab1_ex10}, find the voltage $V$ across resistor R2. You can do this in any way, but remember to explain it in detail. Then simulate the circuit to check the result.

\labfig{lab1_ex10.png}{9cm}{Exercise 10: find the voltage $V$}{fig:lab1_ex10}

\theory{11}

\simulation{
\item Draw the circuit of Figure~\ref{fig:lab1_ex10} in LTspice (components \cmd{res}/\cmd{voltage}/\cmd{current}, ground \cmd{0}) and save it as \file{Lab1/Ex10.asc}.
\item Add the \cmd{.op} analysis (\textit{Simulate $\rightarrow$ Edit Simulation Cmd $\rightarrow$ DC op pnt}), run the simulation and read the node voltages and branch currents in the operating-point window.
\item Label node E as the ground (\texttt{0}) and the other nodes A to D.
\item Be careful with signs: LTspice reports the current of a voltage source as positive when it flows \emph{into} its positive terminal, and the current of a resistor as positive when it flows from its first pin to its second pin.
}

\comparetable{
\cmprow{$V$ (V)}
}

\commentlines{2}

\newpage
\section{Advanced exercises}
\subsection{Advanced Exercise 1: maximum power transfer}
\status{Advanced / Supplementary (maximum power transfer theorem with parametric sweep \cmd{.step} and \cmd{.meas}).}
In the circuit of Figure~\ref{fig:lab1_adv01}, a 12\,V source with a 100\,$\Omega$ series resistor R1 feeds node $x$; a 220\,$\Omega$ resistor R2 connects node $x$ to ground, and a load resistor $R_L$ is also connected between node $x$ and ground.
\labfig{lab1_adv01.png}{9cm}{Advanced Exercise 1: circuit with a variable load $R_L$}{fig:lab1_adv01}
\subsubsection{Theory calculation}\vspace{-3pt}
\notesblock
\begin{enumerate}[label=\alph*.]
  \item Find the Thevenin equivalent (voltage $V_{th}$ and resistance $R_{th}$) seen by $R_L$.\\ \dotfill
  \item Express the power $P_L$ dissipated in $R_L$ as a function of $R_L$, and determine the value of $R_L$ that maximises $P_L$ and the maximum power.\\ \dotfill
\end{enumerate}
\dotlines{4}
\simulation{
\item Draw the circuit of Figure~\ref{fig:lab1_adv01} in LTspice and save it as \file{Lab1/Ex11.asc}.
\item Define the load as a parameter: \cmd{.param RL=100}, and enter \cmd{\{RL\}} as the value of the load resistor.
\item Add \cmd{.step param RL 10 200 10} to simulate $R_L$ from 10 to 200\,$\Omega$ in steps of 10\,$\Omega$.
\item To evaluate one measurement per step, use a one-point DC sweep of the source: \cmd{.dc V1 0 12 12} (the sweep goes from 0\,V to 12\,V in one 12\,V step).
\item Measure the load power in every step with \cmd{.meas DC PL FIND V(x)*I(RL) AT=12}; the results appear in the SPICE Error Log. If your version rejects this expression, measure \cmd{V(x)} and \cmd{I(RL)} separately and compute the product yourself.
\item Fill in the table below (use the steps closest to the listed values) and find the load for which $P_L$ is maximum.
}
\begin{center}\renewcommand{\arraystretch}{1.7}
\begin{tabular}{|c|c|c|c|}\hline
\textbf{$R_L$ ($\Omega$)} & \textbf{$V_L$ (V)} & \textbf{$I_L$ (mA)} & \textbf{$P_L$ (mW)}\\ \hline
20 & & & \\ \hline 40 & & & \\ \hline 60 & & & \\ \hline 80 & & & \\ \hline 100 & & & \\ \hline 150 & & & \\ \hline 200 & & & \\ \hline
\end{tabular}\end{center}
\commentlines{2}

\newpage
\subsection{Advanced Exercise 2: superposition with \texttt{.step}}
\status{Advanced / Supplementary (superposition theorem verification using \cmd{.step} and \cmd{.dc}).}
In the circuit of Figure~\ref{fig:lab1_adv02}, two sources feed node $x$: $V_1=12$\,V through R1\,=\,1\,k$\Omega$ and $V_2=6$\,V through R2\,=\,2.2\,k$\Omega$; resistor R3\,=\,3.3\,k$\Omega$ connects node $x$ to ground.
\labfig{lab1_adv02.png}{9cm}{Advanced Exercise 2: two-source circuit}{fig:lab1_adv02}
\subsubsection{Theory calculation}\vspace{-3pt}
\notesblock
Find $V_x$ (a) with only $V_1$ active, (b) with only $V_2$ active, and (c) with both sources active. Show that (c) equals the sum of (a) and (b).
\dotlines{6}
\simulation{
\item Draw the circuit of Figure~\ref{fig:lab1_adv02} in LTspice and save it as \file{Lab1/Ex12.asc}.
\item Make the value of $V_1$ depend on a parameter: set it to \cmd{\{12*k1\}} and place \cmd{.param k1=1}.
\item Add \cmd{.step param k1 list 0 1} (source $V_1$ off, then on) and a one-point sweep of $V_2$: \cmd{.dc V2 0 6 6} (source $V_2$ off, then on).
\item Measure $V_x$ with \cmd{.meas DC Vx\_V2off FIND V(x) AT=0} and \cmd{.meas DC Vx\_V2on FIND V(x) AT=6}; read the four values (one per combination) in the SPICE Error Log.
}
\begin{center}\renewcommand{\arraystretch}{1.8}
\begin{tabular}{|p{5.2cm}|c|c|c|}\hline
\textbf{Case} & \textbf{Theory (V)} & \textbf{LTspice (V)} & \textbf{Error (\%)}\\ \hline
Only $V_1$ ($V_2$ off) & & & \\ \hline
Only $V_2$ ($V_1$ off) & & & \\ \hline
Both sources on & & & \\ \hline
Sum of the two single-source cases & & & \\ \hline
\end{tabular}\end{center}
\commentlines{2}

\section*{Table of standard resistor values}
\addcontentsline{toc}{section}{Table of standard resistor values (for Exercise 5)}
Standard values for resistors with 10\,\% tolerance (E12 series, in ohms; multiply by powers of ten):
\begin{center}\renewcommand{\arraystretch}{1.4}
\begin{tabular}{|*{12}{c|}}\hline
10 & 12 & 15 & 18 & 22 & 27 & 33 & 39 & 47 & 56 & 68 & 82\\ \hline
\end{tabular}\end{center}
